% ============================================================
\chapter{Points Fixes Al\'eatoires}
\label{ch:aleatoires}
% ============================================================

\begin{intuition}
La th\'eorie des points fixes al\'eatoires \'etend les r\'esultats classiques au cadre
stochastique~: on consid\`ere des op\'erateurs $T(\omega, \cdot)$ qui d\'ependent d'un
param\`etre al\'eatoire $\omega \in \Omega$, et l'on cherche un point fixe $x^*(\omega)$
qui soit lui-m\^eme mesurable en $\omega$. Les difficult\'es principales sont de nature
mesurable~: il faut s'assurer que les constructions (it\'erations, passages \`a la limite)
pr\'eservent la mesurabilit\'e. Ces r\'esultats trouvent des applications en analyse
stochastique, en \'economie math\'ematique et en th\'eorie des jeux stochastiques.
\end{intuition}

% ------------------------------------------------------------
\section{Op\'erateurs al\'eatoires}
% ------------------------------------------------------------

\begin{definition}[Espace mesurable s\'eparable]
\index{espace mesurable s\'eparable}
Un espace topologique $(X, \tau)$ est \emph{s\'eparable} s'il contient un sous-ensemble
d\'enombrable dense. On munit $X$ de sa tribu bor\'elienne
$\mathcal{B}(X) = \sigma(\tau)$.
\end{definition}

\begin{definition}[Op\'erateur al\'eatoire]
\index{op\'erateur al\'eatoire}
Soient $(\Omega, \mathcal{A}, \PP)$ un espace de probabilit\'e et $(X, d)$ un espace
m\'etrique s\'eparable. Un \emph{op\'erateur al\'eatoire} est une application
$T\colon \Omega \times X \to X$ telle que~:
\begin{enumerate}
  \item pour tout $x \in X$, l'application $\omega \mapsto T(\omega, x)$ est
    $(\mathcal{A}, \mathcal{B}(X))$-mesurable~;
  \item pour tout $\omega \in \Omega$ (ou $\PP$-presque tout $\omega$), l'application
    $x \mapsto T(\omega, x)$ satisfait certaines propri\'et\'es (continuit\'e, contraction,
    etc.).
\end{enumerate}
\end{definition}

\begin{definition}[Point fixe al\'eatoire]
\index{point fixe al\'eatoire}
Un \emph{point fixe al\'eatoire} de $T$ est une variable al\'eatoire
$\xi\colon \Omega \to X$ (mesurable) telle que
\[
  T(\omega, \xi(\omega)) = \xi(\omega) \quad \text{pour } \PP\text{-presque tout } \omega \in \Omega.
\]
\end{definition}

\begin{remarque}
L'existence d'un point fixe pour chaque $\omega$ (par un th\'eor\`eme d\'eterministe) ne
garantit pas l'existence d'un point fixe al\'eatoire~: la mesurabilit\'e de
$\omega \mapsto x^*(\omega)$ est un probl\`eme non trivial.
\end{remarque}

% ------------------------------------------------------------
\section{Th\'eor\`eme de s\'election mesurable}
% ------------------------------------------------------------

\begin{definition}[Correspondance mesurable]
\index{correspondance mesurable}
Une correspondance $F\colon \Omega \rightrightarrows X$ est \emph{mesurable} si pour tout
ouvert $U \subseteq X$, l'ensemble $\{\omega : F(\omega) \cap U \neq \emptyset\}$
appartient \`a $\mathcal{A}$.
\end{definition}

\begin{theoreme}[S\'election mesurable de Kuratowski--Ryll-Nardzewski]
\label{thm:selection-mesurable}
\index{th\'eor\`eme de Kuratowski--Ryll-Nardzewski}
Soit $(\Omega, \mathcal{A})$ un espace mesurable et $(X, d)$ un espace m\'etrique complet
s\'eparable. Si $F\colon \Omega \rightrightarrows X$ est une correspondance mesurable
\`a valeurs ferm\'ees non vides, alors $F$ admet une \emph{s\'election mesurable}~:
il existe $f\colon \Omega \to X$ mesurable telle que $f(\omega) \in F(\omega)$ pour
tout $\omega$.
\end{theoreme}

\begin{proof}[Id\'ee de la preuve]
On construit une suite $(f_n)$ de s\'elections $\varepsilon_n$-approch\'ees (avec
$\varepsilon_n \to 0$) par r\'ecurrence. Soit $(x_k)_{k \geq 1}$ une suite dense dans $X$.
Pour chaque $n$, on partitionne $\Omega$ en ensembles mesurables
$A_k^n = \{\omega : F(\omega) \cap B(x_k, \varepsilon_n) \neq \emptyset\} \setminus
\bigcup_{j<k} A_j^n$ et on pose $f_n(\omega) = x_k$ pour $\omega \in A_k^n$. La suite
$(f_n)$ converge uniform\'ement vers une s\'election mesurable $f$.
\end{proof}

% ------------------------------------------------------------
\section{Th\'eor\`eme de point fixe al\'eatoire de Banach}
% ------------------------------------------------------------

\begin{theoreme}[Point fixe al\'eatoire contractant]
\label{thm:pf-aleatoire-banach}
Soit $(\Omega, \mathcal{A}, \PP)$ un espace de probabilit\'e, $(X, d)$ un espace m\'etrique
complet s\'eparable, et $T\colon \Omega \times X \to X$ un op\'erateur al\'eatoire tel que~:
\begin{enumerate}
  \item pour tout $x \in X$, $\omega \mapsto T(\omega, x)$ est mesurable~;
  \item pour $\PP$-presque tout $\omega$, $x \mapsto T(\omega, x)$ est une contraction
    de constante $k(\omega) < 1$.
\end{enumerate}
Alors $T$ admet un unique point fixe al\'eatoire.
\end{theoreme}

\begin{proof}
Pour chaque $\omega$, le th\'eor\`eme de Banach d\'eterministe fournit un unique point
fixe $\xi(\omega) \in X$. Il reste \`a montrer que $\xi$ est mesurable.

Fixons $x_0 \in X$ et d\'efinissons la suite d'it\'er\'ees~:
\[
  \xi_0(\omega) = x_0, \qquad \xi_{n+1}(\omega) = T(\omega, \xi_n(\omega)).
\]
Par hypoth\`ese, $\xi_0$ est mesurable (constante). Si $\xi_n$ est mesurable, alors
$\omega \mapsto T(\omega, \xi_n(\omega))$ est mesurable comme composition de fonctions
mesurables (th\'eor\`eme de Carath\'eodory). Donc $\xi_{n+1}$ est mesurable.

Pour presque tout $\omega$, $\xi_n(\omega) \to \xi(\omega)$ (convergence du th\'eor\`eme
de Banach). Comme limite ponctuelle de fonctions mesurables, $\xi$ est mesurable.
\end{proof}

\begin{attention}
L'hypoth\`ese de s\'eparabilit\'e de $X$ est essentielle pour garantir la mesurabilit\'e
conjointe et la mesurabilit\'e de la limite. Sans cette hypoth\`ese, le r\'esultat est
faux en g\'en\'eral.
\end{attention}

% ------------------------------------------------------------
\section{Point fixe al\'eatoire de Schauder}
% ------------------------------------------------------------

\begin{theoreme}[Point fixe al\'eatoire de Schauder]
\label{thm:pf-aleatoire-schauder}
Soit $(\Omega, \mathcal{A})$ un espace mesurable, $C$ un sous-ensemble convexe ferm\'e
born\'e d'un espace de Banach s\'eparable $E$, et $T\colon \Omega \times C \to C$ un
op\'erateur al\'eatoire tel que~:
\begin{enumerate}
  \item pour tout $x \in C$, $\omega \mapsto T(\omega, x)$ est mesurable~;
  \item pour tout $\omega$, $x \mapsto T(\omega, x)$ est continue~;
  \item pour tout $\omega$, $T(\omega, \cdot)$ est compacte (envoie les born\'es sur des
    ensembles relativement compacts).
\end{enumerate}
Alors $T$ admet un point fixe al\'eatoire.
\end{theoreme}

\begin{proof}[Id\'ee de la preuve]
Par le th\'eor\`eme de Schauder d\'eterministe, pour chaque $\omega$, l'ensemble des points
fixes $F(\omega) = \{x \in C : T(\omega, x) = x\}$ est non vide et ferm\'e. On montre que
la correspondance $\omega \mapsto F(\omega)$ est mesurable. Le th\'eor\`eme de s\'election
mesurable de Kuratowski--Ryll-Nardzewski fournit alors une s\'election mesurable, qui est
le point fixe al\'eatoire cherch\'e.
\end{proof}

% ------------------------------------------------------------
\section{Applications en analyse stochastique}
% ------------------------------------------------------------

\begin{proposition}[\'Equations diff\'erentielles al\'eatoires]
Consid\'erons l'\'equation diff\'erentielle al\'eatoire
\[
  y'(t, \omega) = f(\omega, t, y(t, \omega)), \qquad y(t_0, \omega) = Y_0(\omega)
\]
o\`u $Y_0\colon \Omega \to \R^n$ est une variable al\'eatoire et $f$ v\'erifie~:
\begin{itemize}
  \item $f(\omega, \cdot, \cdot)$ est continue en $(t,y)$ pour tout $\omega$~;
  \item $f(\cdot, t, y)$ est mesurable en $\omega$ pour tous $t, y$~;
  \item $f$ est lipschitzienne en $y$ uniform\'ement en $\omega$.
\end{itemize}
Alors il existe une unique solution $y(\cdot, \omega)$ qui est un processus mesurable.
\end{proposition}

\begin{proof}
L'op\'erateur de Picard $T(\omega, u)(t) = Y_0(\omega) + \int_{t_0}^t f(\omega, s, u(s)) ds$
est un op\'erateur al\'eatoire contractant sur $C([t_0-\delta, t_0+\delta], \R^n)$ muni
de la norme de Bielecki. Le Th\'eor\`eme~\ref{thm:pf-aleatoire-banach} s'applique.
\end{proof}

\begin{proposition}[\'Equations int\'egrales stochastiques]
L'\'equation int\'egrale al\'eatoire de Fredholm
\[
  u(\omega, t) = f(\omega, t) + \lambda \int_a^b K(\omega, t, s) u(\omega, s) \, ds
\]
admet un unique point fixe al\'eatoire d\`es que $\abs{\lambda} < 1/((b-a) \sup_\omega \norm{K(\omega, \cdot, \cdot)}_\infty)$.
\end{proposition}

% ------------------------------------------------------------
\section{Th\'eor\`eme d'Itoh}
% ------------------------------------------------------------

\begin{theoreme}[Itoh]
\label{thm:itoh}
\index{th\'eor\`eme d'Itoh}
Soit $(\Omega, \mathcal{A})$ un espace mesurable complet, $C$ un sous-ensemble ferm\'e
s\'eparable d'un espace de Banach, et $T\colon \Omega \times C \to C$ un op\'erateur
al\'eatoire non expansif ($d(T(\omega,x), T(\omega,y)) \leq d(x,y)$) tel que pour
tout $\omega$, $T(\omega, C)$ est contenu dans un compact de $C$. Alors $T$ admet
un point fixe al\'eatoire.
\end{theoreme}

\begin{formules}[title=\textbf{R\'ecapitulatif des th\'eor\`emes de points fixes al\'eatoires}]
\begin{center}
\renewcommand{\arraystretch}{1.3}
\begin{tabular}{|l|l|l|}
\hline
\textbf{Th\'eor\`eme} & \textbf{Hypoth\`ese sur $T(\omega, \cdot)$} & \textbf{Unicit\'e} \\
\hline
PF al\'eatoire de Banach & Contraction & Oui \\
PF al\'eatoire de Schauder & Continue, compacte & Non \\
Itoh & Non expansif, image compacte & Non \\
\hline
\end{tabular}
\end{center}
\end{formules}

% ------------------------------------------------------------
\section{Exercices}
% ------------------------------------------------------------

\begin{exercice}
Soit $\Omega = [0,1]$ muni de la mesure de Lebesgue et $X = \R$. D\'efinir
$T(\omega, x) = \omega x / 2$. V\'erifier que $T$ est un op\'erateur al\'eatoire
contractant et d\'eterminer explicitement le point fixe al\'eatoire.
\end{exercice}

\begin{exercice}
Soit $T(\omega, x) = a(\omega) x + b(\omega)$ o\`u $a, b\colon \Omega \to \R$ sont
mesurables avec $\abs{a(\omega)} < 1$ p.s. Montrer que le point fixe al\'eatoire est
$\xi(\omega) = b(\omega) / (1 - a(\omega))$. V\'erifier sa mesurabilit\'e.
\end{exercice}

\begin{exercice}
Soit $(X, d)$ un espace m\'etrique compact et $T\colon \Omega \times X \to X$ un
op\'erateur al\'eatoire continu. Montrer que la correspondance
$F(\omega) = \{x : T(\omega, x) = x\}$ est mesurable.
\emph{Indication~: montrer que $\{\omega : F(\omega) \cap U \neq \emptyset\}$ est
mesurable pour tout ouvert $U$.}
\end{exercice}

\begin{exercice}
Consid\'erons l'EDO al\'eatoire $y' = A(\omega) y$, $y(0) = 1$, o\`u $A(\omega)$ est une
variable al\'eatoire r\'eelle. Reformuler comme probl\`eme de point fixe et appliquer le
th\'eor\`eme de point fixe al\'eatoire contractant.
\end{exercice}

\begin{exercice}
Donner un exemple d'op\'erateur al\'eatoire $T\colon \Omega \times [0,1] \to [0,1]$
continu en $x$ et mesurable en $\omega$, qui admet un point fixe pour chaque $\omega$,
mais pour lequel il n'existe aucune s\'election mesurable de points fixes.
\emph{Indication~: utiliser un espace $(\Omega, \mathcal{A})$ non complet.}
\end{exercice}

\begin{exercice}[It\'erations al\'eatoires]
\index{syst\`eme de fonctions it\'er\'ees}
Soit $(T_1, \ldots, T_k)$ une famille de contractions de $\R^n$ avec des probabilit\'es
$(p_1, \ldots, p_k)$. On d\'efinit $x_{n+1} = T_{\sigma_n}(x_n)$ o\`u $(\sigma_n)$ est
i.i.d.\ de loi $(p_i)$. Montrer que la suite $(x_n)$ converge p.s.\ vers une mesure
invariante (attracteur al\'eatoire), ind\'ependamment du point initial $x_0$.
\end{exercice}
